\section{Method}
\label{sec:method}

The method was pre-registered in a versioned protocol document before any corpus measurement, and the
commit timestamp is its evidence. This section reports what was actually executed, including where it
departs from the pre-registered plan; departures are enumerated in Section~\ref{sec:threats} and in
the replication package.

\subsection{Sampling frame}

A defensible frame must be reproducible, must not be ranked by popularity, and must not depend on
credentials the study does not have. The GitHub code-search API was unavailable (no access token;
the unauthenticated quota is 60 requests per hour) and is in any case capped at approximately 1{,}000
results and ranked by popularity --- characteristics that would have made ``the top results'' both
unrepresentative and unstable.

We therefore define the frame as follows:

\begin{quote}
\textbf{Frame.} All public GitHub repositories that had adopted at least one of twelve declared
BDD-related topic tags --- \texttt{cucumber}, \texttt{bdd}, \texttt{gherkin}, \texttt{behave},
\texttt{cucumber-jvm}, \texttt{cucumber-js}, \texttt{pytest-bdd}, \texttt{specflow},
\texttt{behavior-driven-development}, \texttt{bdd-framework}, \texttt{acceptance-testing},
\texttt{living-documentation} --- as observed on 25 November 2026, enumerated across the language
filters \emph{any}, \emph{java}, \emph{javascript}, \emph{typescript} and \emph{python}, six pages per
combination.
\end{quote}

Enumeration produced \Nframe{} unique repositories from 252 topic pages, with the SHA-256 of the
concatenated page payload recorded so that the snapshot is auditable. The sample is a
\emph{hash-ordered permutation} of the frame with fixed seed \Seed{}: repositories are ordered by
$\mathrm{SHA256}(\texttt{"<seed>:<repo>"})$, which is reproducible in any language and independent of
the random-number generator of any particular runtime.

We state the frame's bias explicitly rather than concealing it: repositories that never adopted a BDD
topic tag are out of frame, and topic tagging correlates with project self-presentation. The frame is
therefore a sample of \emph{self-identified} BDD projects, not of all projects using Cucumber.

\subsection{Eligibility}

A repository enters the corpus only if all eight criteria hold, each mechanically checkable:

\begin{enumerate}
\item at least one \texttt{.feature} file;
\item at least 20 scenarios, counting \texttt{Scenario} and \texttt{Scenario Outline} headings (and
      their Portuguese equivalents);
\item at least one step-definition file;
\item at least one build descriptor (\texttt{pom.xml}, \texttt{build.gradle}, \texttt{package.json},
      \texttt{pyproject.toml}, \texttt{setup.py}, \ldots);
\item a language in scope: Java, JavaScript, TypeScript or Python;
\item a detectable permissive licence;
\item \emph{application shape}: at least ten non-test source files, because the construct requires a
      domain layer to mutate;
\item a BDD runner named in a build manifest or dependency file.
\end{enumerate}

Criteria 7 and 8 were added after the pilot (P1.5) revealed that the naive rule --- features, steps,
build, licence --- admitted reporting agents, test utilities and Maven plugins but almost no
applications with business rules. This is a change to the frame definition made \emph{before} the
gate closed, with the supporting pilot data, and it is disclosed here and in the replication package.

\subsection{Funnel and attrition}

Every candidate is recorded in exactly one funnel state with a categorised reason, and this log
\emph{is} the attrition dataset:

\begin{quote}
\texttt{eligible} $\rightarrow$ \texttt{bdd\_passes} \quad|\quad \texttt{bdd\_fails} \quad|\quad
\texttt{no\_bdd\_execution} \quad|\quad \texttt{build\_fails} \quad|\quad \texttt{timeout}
\end{quote}

with exclusion reasons drawn from a fixed vocabulary: \texttt{no\_feature\_files},
\texttt{too\_few\_scenarios}, \texttt{no\_step\_definitions}, \texttt{no\_build\_config},
\texttt{language\_out\_of\_scope}, \texttt{licence\_not\_permissive},
\texttt{no\_bdd\_runner\_dependency}, \texttt{no\_application\_shape}, and, for build attempts,
\texttt{dependency\_rot}, \texttt{tool\_incompatibility}, \texttt{missing\_service\_dependency},
\texttt{needs\_browser}, \texttt{needs\_database}, \texttt{compile\_error}, \texttt{timeout},
\texttt{other}.

One distinction is load-bearing and is enforced in the harness rather than left to the analyst.
\texttt{no\_bdd\_execution} denotes a repository whose tests pass but where no evidence exists that
the Gherkin layer ran at all; \texttt{bdd\_fails} denotes a repository whose tests ran and failed.
Reporting the first as the second would inflate apparent build attrition, and reporting either as
\texttt{bdd\_passes} would admit a suite that never executed a scenario. Because a green
\texttt{mvn test} proves nothing about Gherkin, the harness requires positive evidence of Cucumber
execution in the output before a repository is accepted.

\subsection{Outcome: oracle strength}

For repository $r$, let $M$ be the generated domain-layer mutants, $K$ the killed mutants and
$\mathit{Cov}$ those executed by the run in question. Statuses are taken from the tool's
\emph{per-mutant} records, never from its summary line, because the tools disagree about what their
headline means.

\begin{align}
\msall(\textit{mode}) &= \frac{|K| + |T|}{|K| + |S| + |T| + |\mathit{NC}|} \label{eq:msall}\\
\mscov(\textit{mode}) &= \frac{|K| + |T|}{|K| + |S| + |T|} \label{eq:mscov}\\
\mathit{coverage}(\textit{mode}) &= \frac{|K| + |S| + |T|}{|K| + |S| + |T| + |\mathit{NC}|} \label{eq:cov}
\end{align}

where $S$ denotes surviving, $T$ timed-out and $\mathit{NC}$ uncovered mutants. Timeouts count as
detected in the primary formula, on the reasoning that a mutant which prevents termination was
distinguished from the original; a sensitivity removing them is also reported. Invalid mutants
(compile errors, non-viable mutations) are excluded from every denominator \emph{and their count is
reported}. The primary outcome is $\mscov(\bdd)$.

Coverage conditioning is mandatory and the following rule is binding on every table:

\begin{quote}
\label{quot:coverage}
A low $\msall$ produced by low coverage says that the BDD suite does not execute that code --- ``not
covered'' --- and is \emph{not} evidence of drift, nor of a weak oracle over the behaviour the suite
does exercise. Only $\mscov$ speaks to the discriminating power of the suite over the domain
behaviour it covers.
\end{quote}

\textbf{Differential kill.} Let $\mathcal{S}$ be the mutants executed in \emph{both} modes. Formally
the protocol asserts $\mathit{executed}(\bdd) \subseteq \mathit{executed}(\full)$, so
$\mathcal{S}$ is the BDD-covered set and every mutant in it is also executed by the full suite. On
that set we report

\begin{align}
\dk &= \mscov(\full,\mathcal{S}) - \mscov(\bdd,\mathcal{S}) \label{eq:dk}\\
\mathit{DK}_{\mathit{count}} &= \frac{|K_{\full}(\mathcal{S}) \setminus K_{\bdd}(\mathcal{S})|}{|\mathcal{S}|}
\qquad
\mathit{BDDshare} = \frac{|K_{\bdd}(\mathcal{S})|}{|K_{\full}(\mathcal{S})|} \label{eq:dkshare}
\end{align}

$\dk$ answers \emph{what the specification layer adds beyond the repository's other tests};
$\mathit{BDDshare}$ answers \emph{how much of the full suite's oracle is the specification layer}.
Both are reported, with $\dk$ as the headline because it is a difference of scores on an identical
mutant set.

\subsection{Drift indicators}

Seven indicators are defined so that higher always means more drift-like. None is sign-flipped after
measurement. Denominators use \emph{active} scenarios, i.e.\ all scenarios minus those already
flagged by the skip indicator.

\textbf{D1 --- \texttt{skip\_ratio} (primary).} A scenario is skipped if a whole-tag match for
\texttt{@skip}, \texttt{@ignored}, \texttt{@wip}, \texttt{@pending}, \texttt{@todo},
\texttt{@manual} or equivalents appears on the scenario, its \texttt{Examples:} block, or its feature:
\begin{equation}
\texttt{skip\_ratio} = \frac{|\{s \in S_r : s \text{ is skipped}\}|}{|S_r|}.
\end{equation}

\textbf{D2 --- \texttt{assertion\_weakness\_ratio} (primary).} The unit is one step definition. A step
definition is \texttt{STRONG\_DIRECT} if it carries an assertion token and matches no weakening
pattern; \texttt{STRONG\_DELEGATED} if it carries no token but calls an intra-repository method that
is strong, resolved to depth $\le 2$; \texttt{WEAK} if it carries a token and matches a weakening
pattern; and \texttt{NONE} if it has neither token nor resolvable callee. Weakening patterns are
tautology (identical operands), constant truth, unreachable assertion (after an unconditional return,
or in a literal-conditioned branch), swallowed failure (a \texttt{catch} that neither rethrows nor
fails), and self-comparison:
\begin{equation}
\texttt{assertion\_weakness} = \frac{|\texttt{WEAK}| + |\texttt{NONE}|}
{|\texttt{STRONG\_DIRECT}| + |\texttt{STRONG\_DELEGATED}| + |\texttt{WEAK}| + |\texttt{NONE}|}.
\end{equation}

\textbf{D3 --- \texttt{stale\_spec\_ratio}.} For an active scenario, let $t_s$ be the date of its
most recent touching commit (from \texttt{git blame} over its line range) and $n_s$ the number of
commits since $t_s$ touching the domain layer. Then
$\texttt{stale} = |\{s : \mathrm{age}_s \ge 180 \wedge n_s \ge 5\}| / |A_r|$, with the thresholds 180
days and 5 commits fixed before measurement.

\textbf{D4 --- \texttt{redundancy\_cluster\_ratio}.} Active scenarios are fingerprinted by the set of
their normalised step texts (parameters replaced by placeholders), an edge joins two scenarios at
Jaccard $\ge 0.8$, and $\texttt{redundancy} = 1 - (\text{connected components})/|A_r|$.

\textbf{D5 --- \texttt{glue\_duplication\_ratio}.} Distinct normalised step invocations divided by
total step invocations across active scenarios. Scripted one-to-one glue approaches 1; parameterised,
reused steps fall well below.

\textbf{D6 --- \texttt{literal\_bearing\_stepdef\_ratio}.} The share of resolved step definitions
containing at least one counted literal: numeric literals other than $0$, $1$ and $-1$, and string
literals of length $\ge 2$. Excluded, because they are mechanism rather than business constants:
booleans and nulls, enum constants, the step expression itself, assertion messages, log and format
strings, and UI constants (selectors, URLs, XPaths).

\textbf{D7 --- \texttt{mirror\_stepdef\_ratio}.} A \emph{mirror} step re-derives the expected value by
repeating the production computation, so that its oracle moves with the mutant. The rater-defined
form requires that the asserted expected value be computed (not a literal), that the computation
reference a domain-layer call in the same body, and that the actual value also come from a
domain-layer call, with the rater citing the domain method inspected. The \emph{automated census}
computes the loose conjunction of the first three conditions; its precision is \emph{not} validated
without human raters and is reported as such.

\subsection{Mutation protocol}

Mutants are restricted to operators touching decision logic: conditionals boundary, negated
conditionals, arithmetic, increments, negated values, and return-value changes. The exact operator set
is stated per tool in the replication package, and no operator set was altered after seeing results.
Mutant sampling, where required by runtime, uses the fixed study seed under a per-repository budget.
Timeouts, uncovered mutants and equivalent mutants are counted and reported; none is silently dropped.
The domain layer is delimited by the three-tier rule of protocol step 1, and the tier used is recorded
per repository.

\subsection{Statistics}

The statistical plan was fixed before measurement:

\begin{itemize}
\item \textbf{Unit of analysis} is the repository. Mutants, scenarios and steps are never treated as
      independent observations.
\item \textbf{Exactly two primary tests}: D1 and D2 against $\mscov(\bdd)$, Holm-corrected across the
      pair. All other indicators form one exploratory family, Holm-corrected within that family and
      labelled exploratory in every table.
\item \textbf{No pooled correlation as the headline.} Associations are computed \emph{within}
      language/tool strata first, reported as Spearman $\rho$ with bootstrap 95\,\% confidence
      intervals (2{,}000 resamples, seed \Seed{}). A pooled figure, where given, is the average of
      within-stratum associations with a cluster bootstrap by stratum, and is described as exactly
      that.
\item \textbf{Coverage is a covariate} in every model and appears in every table.
\item \textbf{Power is reported as a detectable-effect floor.} On the Fisher-$z$ scale with
      $\mathrm{SE}(z) = 1.030/\sqrt{n-3}$, the smallest detectable $|\rho|$ at two-sided
      $\alpha = 0.05$ and power $0.80$ is
      \begin{equation}
      |\rho|_{\min} = \tanh\!\left((z_{1-\alpha/2} + z_{1-\beta})\cdot\frac{1.030}{\sqrt{n-3}}\right),
      \end{equation}
      giving $0.80$ at $n = 10$, $0.60$ at $n = 20$, $0.50$ at $n = 30$ and $0.44$ at $n = 40$; under
      Holm ($\alpha = 0.025$ for the smaller of two $p$-values) the floor at $n = 30$ rises to $0.54$.
      An association below the floor at the achieved $n$ is reported as \emph{not detectable at this
      $n$}, which is not the same statement as \emph{absent}.
\item \textbf{Longitudinal analysis} uses the Theil--Sen slope of each indicator against commit time,
      reported with the median slope, the interquartile range, and the proportion of repositories with
      a positive slope with an exact binomial interval. A two-snapshot change in oracle strength is
      reported only on the common mutant set and only if that set covers at least half of the HEAD
      mutants; otherwise the repository is flagged \texttt{low\_overlap}. We call such a change a
      change between two snapshots, not a trend.
\end{itemize}

\subsection{Reproducibility}

Every measurement is recorded in an append-only run log with the command, UTC timestamp, the git
commit of the harness, host, seed and exact tool versions. Raw per-mutant statuses, the screening
tables and the analysis scripts are part of the replication package. The bibliography is generated
from a Crossref-verified table, and every DOI is re-resolved by an automated audit before the
manuscript is built.

